\begin{definition}
    Оператор $A\!: E_1 \to E_2$ называется \textit{обратимым}, если
    \[\forall y \in \text{Im}A \;\; \exists!x \in E_1:\: Ax=y.\]
    $A^{-1}\!:\: E_2 \to E_1$ называется \textit{обратным оператором}, если $\left(\forall x \in E_1\: A^{-1}Ax=x \right) \wedge \left(\forall y \in E_2\: AA^{-1}y=y\right)$.
\end{definition}

\begin{definition}
    Оператор $R$ называется \textit{правым обратным} к $A$, если $\forall y \in \im A \left(ARy=y \right)$, т.е. $AR=I|_{\im A}$

    Оператор $L$~--- левый обратный оператор определяется аналогичным образом. 
    Очевидно, что $R$ отвечает за существование решения, а $L$ за его единственность.
    Тогда понятно что $A^{-1} = L = R$. Такое введение определяется тем, что в разных ситуациях проще искать разный тип обратного оператора.
\end{definition}

\begin{unnecessary}
    \begin{exercise}
        $A$~--- линейный и обратим, то обратный ему оператор линеен
    \end{exercise}

    \begin{example}
        Не каждый линейный оператор обратим~--- пример тождественный ноль.
    \end{example}
    
    \begin{exercise}
    Пусть $A$~--- ограниченный оператор и у него существует обратный $A^{-1}$. Верно ли, что $A^{-1}$~--- ограниченный?
    \end{exercise}
    \begin{answer}
    Нет, это неверно, пример~--- оператор Вольтера на $C[0;1]$:
    \begin{align*}
        (Vf)(x) &=\int_{t=0}^x f(t)dt,\\
        \im \: V &= \{ g \in C^1, g(0)=0 \}.
    \end{align*}
    Обратный оператор~--- оператор дифференцирования, не является ограниченным
    \end{answer}
\end{unnecessary}

\begin{theorem}[8.4, теорема Банаха об обратном операторе]\label{theorem8.4_pre}\label{theorem8.4}
    Пусть $E_1$, $E_2$~--- банаховы пространства, $A \in \mathcal{L}(E_1, E_2)$, причём $A$~--- биекция. Тогда $A^{-1}$~--- линейный ограниченный оператор.
\end{theorem}

\begin{proof}
    В общем случае: б/д. Для гильбертовых пространств дальше
\end{proof}

\begin{definition}[Усиление условия $\ker A = \{0\}$]
    Назовём ЛНО $A$ \textit{ограниченным снизу}, если
    \[\exists m > 0\!:\: \forall x \; \|Ax\| \geqslant m \|x\|.\]
\end{definition}

\begin{note}
    Ограниченность снизу влечёт инъективность
\end{note}

\begin{proof}
    Пусть $A$ не инъективен. Рассмотрим $x \in \ker A, x \neq 0$
        \[ 0 = \| Ax \| \geq m \| x \| \neq 0 \]
    Пришли к противоречию.
\end{proof}

\begin{theorem}[8.1]\label{8.1}
    Пусть $A \in \mathcal{L}(E_1, E_2)$. Существует $A^{-1} \in \L(\im A, \text{dom}A) \Leftrightarrow A$~--- ограничен снизу.
\end{theorem}

\begin{proof}
    \begin{itemize}
        \item[$\Leftarrow$]
        Берём $\forall y \in \im\: A,$ для него верно \(\|y\| = \|Ax\| \geqslant m \|x\| = m \|A^{-1}y\|\). $A^{-1}$ существует по определению: инъективность следует из ограниченности снизу, а сюръективность очевидна, так как просто рассматриваем оператор на $\im A$.
        
        Отсюда $\|A^{-1} y\| \leqslant \frac{1}{m} \|y\|$.
        \item[$\Rightarrow$]
        \( \forall x \in E_1 \; \|x\|=\|A^{-1} y\| \leqslant \| A^{-1}\| \cdot \|y\|=\|A^{-1}\| \cdot \|Ax\| \Rightarrow \|Ax\| \geqslant \frac{1}{\|A^{-1}\|}\|x\| \)
    \end{itemize}
\end{proof}
\newpage